We analyze how the learned maps change the hidden-state context passed to the frozen drafter.

\subsection{What changes when the interface already works?}
\label{sec:analysis-lora}

For a LoRA-fine-tuned target, identity maps recover the native DFlash interface. Write $z_0=\sum_i F_i h_i$ for this context and $\delta=\sum_i F_i(W_i-I)h_i$ for the learned correction. We run decoding with $z_0$ plus either the component of $\delta$ parallel to $z_0$ or the component perpendicular to it.

\begin{table}[ht!]
\centering
\caption{LoRA-target direction and geometry controls. Bold marks the best acceptance length for each task, and underline marks the second best. The geometry measurements use 1,024 positions per task and are dimensionless.}
\label{tab:key-interventions}
\begin{minipage}[t]{0.43\linewidth}
\vspace{0pt}
\centering
\scriptsize
\textbf{Direction intervention}\par\smallskip
\begin{tabular}{@{}lrr@{}}
\toprule
Context & GSM8K $\tau$ & NanoCoder $\tau$ \\
\midrule
Native $z_0$ & 4.472 & 5.111 \\
Full $z_0+\delta$ & \textbf{5.988} & \textbf{5.552} \\
Parallel only & 4.466 & 5.114 \\
Perpendicular only & \underline{5.924} & \underline{5.527} \\
\bottomrule
\end{tabular}
\end{minipage}\hfill
\begin{minipage}[t]{0.54\linewidth}
\vspace{0pt}
\centering
\scriptsize
\textbf{LoRA-shift geometry}\par\smallskip
\setlength{\tabcolsep}{2pt}
\begin{tabular}{@{}lrrrr@{}}
\toprule
Task & \shortstack{Correction\\cosine} & \shortstack{Shuffled\\cosine} & \shortstack{Radial\\energy} & \shortstack{RMSNorm\\error} \\
\midrule
GSM8K & $-0.0429$ & $-0.0413$ & 25.22\% & 17.55\% \\
NanoCoder & 0.1248 & 0.1052 & 34.82\% & 14.88\% \\
\bottomrule
\end{tabular}
\end{minipage}
\end{table}

The parallel component returns acceptance length to the native level, while the perpendicular component recovers nearly all of the full map's acceptance length (Table~\ref{tab:key-interventions}). Scaling the whole context by two leaves acceptance length unchanged at 5.988 tokens per verification step on GSM8K and 5.552 on NanoCoder. The useful change is mainly directional, consistent with RMSNorm suppressing common rescaling. Appendix~\ref{app:lora-direction} reports every result.

Do the maps reverse the change caused by LoRA fine-tuning? We run the base and LoRA-fine-tuned targets on the same 1,024 saved token positions. At each position, we compare the map correction with the change in fused context caused by LoRA. Exact cancellation would require the correction to point in the opposite direction. We also shuffle the LoRA changes across positions while leaving the corrections fixed. This tests whether each correction is specifically aligned with the change at the same token. The mean correction cosine is $-0.0429$ on GSM8K and $0.1248$ on NanoCoder, close to the shuffled values of $-0.0413$ and $0.1052$ (Table~\ref{tab:key-interventions}). The maps therefore do not systematically reverse the LoRA-induced change token by token. These results support redirection of an already usable context, not literal recovery of the pre-LoRA context. Appendix~\ref{app:lora-geometry} defines each measurement and gives the full procedure.

Rank-128 corrections retain 95.1--99.6\% of full-map latency reduction on GSM8K, while rank 256 retains 57.8--89.3\% on NanoCoder. Much of the measured benefit can survive moderate-rank truncation, though timing varies (Appendix~\ref{app:low-rank-corrections}).

\FloatBarrier
\subsection{Why does reconstruction help across model sizes?}
\label{sec:analysis-transfer}

We first compare the mapped Qwen3-8B features and fused context with those of Qwen3-4B on the same training text. MSE gives lower normalized context error and higher feature cosine than AUF or decay CE (Table~\ref{tab:transfer-probe-summary}). We then hold the prefixes and saved continuations fixed and change only the maps. On held-out MATH, Code, and Chat positions, MSE gives the lowest token cross-entropy and the longest correct prefix before the first wrong prediction. GSM8K is mixed. Thus, the recovered interface improves prediction on three of four workloads, not just the reconstruction measure. Appendix Tables~\ref{tab:app-context-geometry}--\ref{tab:app-layer-geometry} and~\ref{tab:app-heldout-math}--\ref{tab:app-heldout-chat} give the full results.

\begin{table}[ht!]
\centering
\caption{Interface recovery on 4,096 response positions from 64 NuminaMath-derived training examples. Source cosine is reported as the minimum--maximum across the five feature streams. Error and cosine are dimensionless. Best is bold and second best is underlined.}
\label{tab:transfer-probe-summary}
\scriptsize
\setlength{\tabcolsep}{3pt}
\begin{tabular}{@{}lrrr@{}}
\toprule
Recipe & Peak LR & Error $\downarrow$ & \shortstack{Source cosine\\min--max $\uparrow$} \\
\midrule
MSE & $10^{-3}$ & \textbf{0.136} & \textbf{0.944--0.985} \\
AUF & $10^{-4}$ & 0.940 & 0.121--0.267 \\
Decay CE & $10^{-4}$ & \underline{0.900} & \underline{0.127--0.279} \\
\bottomrule
\end{tabular}
\end{table}

Without retraining, maps halfway between token-trained and MSE maps improve live throughput and acceptance length on the tested MATH and Code workloads (Appendix~\ref{app:midpoint-full}).

\noindent\textbf{Does initialization explain the difference?}
To test whether Xavier initialization explains MSE's advantage, we retrain from a rectangular identity, which copies the coordinates that fit the narrower interface and sets the rest to zero. All other MSE settings stay fixed, so its comparison isolates initialization. MSE acceptance length is nearly unchanged on MATH and GSM8K and slightly lower on Code and Chat, while its throughput falls on all four workloads (Table~\ref{tab:optimization-controls}). AUF and decay CE improve both measurements on every workload but remain below MSE in acceptance length. Appendix Tables~\ref{tab:app-identity-endpoints} and~\ref{tab:app-identity-change} report all endpoints.

\begin{table}[ht!]
\centering
\caption{Initialization results on 128 prompts per task. Each cell shows Xavier $\rightarrow$ rectangular-identity initialization. Throughput is measured in tokens/s, and acceptance length $\tau$ is measured in tokens per verification step. For each dataset and metric, bold marks the best of all six values, and underline marks the second best.}
\label{tab:optimization-controls}
\scriptsize
\setlength{\tabcolsep}{2pt}
\begin{tabular}{@{}clcccc@{}}
\toprule
Metric & Recipe & MATH & GSM8K & Code & Chat \\
\midrule
\multirow{3}{*}{\shortstack{Throughput\\(tokens/s)}} & MSE & \textbf{218.75} $\rightarrow$ \underline{213.70} & \underline{172.35} $\rightarrow$ 166.70 & \textbf{125.98} $\rightarrow$ \underline{119.49} & \textbf{78.31} $\rightarrow$ \underline{76.13} \\
& AUF & 194.22 $\rightarrow$ 206.16 & 162.05 $\rightarrow$ \textbf{175.35} & 58.37 $\rightarrow$ 73.10 & 49.93 $\rightarrow$ 60.77 \\
& Decay CE & 194.69 $\rightarrow$ 203.74 & 163.47 $\rightarrow$ 169.72 & 61.56 $\rightarrow$ 76.02 & 51.31 $\rightarrow$ 61.11 \\
\midrule
\multirow{3}{*}{\shortstack{Acceptance\\length $\tau$\\[-2pt](tokens/step)}} & MSE & \underline{7.973} $\rightarrow$ \textbf{7.983} & \underline{6.204} $\rightarrow$ \textbf{6.217} & \textbf{4.690} $\rightarrow$ \underline{4.569} & \textbf{2.802} $\rightarrow$ \underline{2.747} \\
& AUF & 6.892 $\rightarrow$ 7.326 & 5.677 $\rightarrow$ 6.166 & 2.346 $\rightarrow$ 2.931 & 1.767 $\rightarrow$ 2.163 \\
& Decay CE & 6.947 $\rightarrow$ 7.348 & 5.693 $\rightarrow$ 6.073 & 2.449 $\rightarrow$ 3.042 & 1.812 $\rightarrow$ 2.183 \\
\bottomrule
\end{tabular}
\end{table}
